Mixed-integer programming models used for operational decisions are often produced
by software generators. Their developers must connect numerical feasibility to the
accuracy needed in the application: a small residual in an exported row need not
be small in megawatts, water volume or material mass. Multiplying a constraint and
its right-hand side (RHS) by a positive number preserves the model in exact
arithmetic, but changes both the coefficients presented to the optimizer and the
original-unit meaning of an absolute residual tolerance. Solvers scale internally;
a generator can also use row roles, variable ownership and application units before
export \cite{gurobiExternalScaling,idaesScaling}.

The practical task is to make those accuracy requirements part of model generation
and verification. Coefficient centering alone does not resolve it: scaling a row
down enlarges the original-unit residual admitted by a fixed exported-row tolerance.
Conversely, a factor chosen only for residual accuracy may conflict with coefficient
limits. We connect four steps: declare an original-unit residual budget, select an
admissible factor or report incompatibility, check the written model, and independently
verify the returned point. The application budgets in this study are stated numerical
precision goals, rather than operator standards.
